\documentclass[oneside]{article}
\usepackage[backend=biber,natbib=true,style=authoryear]{biblatex}
\addbibresource{/home/nguyen/1_NQBH/reference/bib.bib}
\usepackage[utf8]{inputenc}
\usepackage{float}
\usepackage{graphicx}
\usepackage[colorlinks=true,linkcolor=blue,urlcolor=red,citecolor=magenta]{hyperref}
\usepackage{amsmath,amssymb,amsthm,mathtools}
\allowdisplaybreaks
\numberwithin{equation}{section}
\newtheorem{assumption}{Assumption}[section]
\newtheorem{lemma}{Lemma}[section]
\newtheorem{corollary}{Corollary}[section]
\newtheorem{definition}{Definition}[section]
\newtheorem{proposition}{Proposition}[section]
\newtheorem{theorem}{Theorem}[section]
\newtheorem{notation}{Notation}[section]
\newtheorem{remark}{Remark}[section]
\newtheorem{example}{Example}[section]
\newtheorem{ques}{Question}[section]
\newtheorem{problem}{Problem}[section]
\newtheorem{conjecture}{Conjecture}[section]
\usepackage[left=0.5in,right=0.5in,top=1.5cm,bottom=1.5cm]{geometry}
\usepackage{fancyhdr}
\pagestyle{fancy}
\fancyhf{}
\lhead{\small \textsc{Sect.} ~\thesection}
\rhead{\small \nouppercase{\leftmark}} % \nouppercase !
\renewcommand{\sectionmark}[1]{\markboth{#1}{}}
\cfoot{\thepage}
\def\labelitemii{$\circ$}

\title{Susan Cain. \textit{Quiet: The Power of Introverts in a World That Can't Stop Talking}}
\author{Collector: Nguyen Quan Ba Hong}
\date{\today}

\begin{document}
\maketitle
\setcounter{secnumdepth}{6}
\setcounter{tocdepth}{6}
\tableofcontents

%------------------------------------------------------------------------------%

\section*{More advance noise for \textit{quiet}}
\begin{quotation}
    ``\textbf{\textit{An intriguing and potentially life-altering examination of the human psyche that is sure to benefit both introverts and extroverts alike}.}'' - \textit{Kirkus Reviews} (starred review)

    ``\textit{Gentle is powerful $\ldots$ Solitude is socially productive $\ldots$ These important counterintuitive ideas are among the many reasons to \textbf{take \emph{Quiet} to a quiet corner and absorb its brilliant, thought-provoking message.}}'' - Rosabeth Moss Kanter, professor at Harvard Business School, author of \textit{Confidence} and \textit{SuperCorp}
\end{quotation}

\begin{quotation}
    ``\textit{\textbf{An informative, well-researched book} on the power of quietness and the virtues of having a rich inner life. It dispels the myth that you have to be extroverted to be happy and successful.}'' - Judith Orloff, M.D., author of \textit{Emotional Freedom}
    
    ``\textit{If this \textbf{engaging and beautifully written} book, Susan Cain makes a powerful case for the wisdom of introspection. She also warns us ably about the downside to our culture's noisiness, including all that it risks drowning out. Above the din, \textbf{Susan's own voice remains a compelling presence-thoughtful, generous, calm, and eloquent. \emph{Quiet} deserves a very large readership.}}'' - Christopher Lane, author of \textit{Shyness: How Normal Behavior Became a Sickness}
    
    ``\textit{Susan Cain's quest to understand introversion, \textbf{a beautifully wrought journey from the lab bench to the motivational speaker's hall}, offers convincing evidence for valuing substance over style, steak over sizzle, and qualities that are, in America, often derided. \textbf{This book is brilliant, profound, full of feeling and brimming with insights.}}'' - Sheri Fink, M.D., author of \textit{War Hospital}
    
    ``\textit{\textbf{Brilliant, illuminating, empowering!} \emph{Quiet} gives not only a voice, but a path to homecoming for so many who've walked through the better part of their lives thinking the way they engage with the world is something in need of fixing.}'' - Jonathan Fields, author of \textit{Uncertainty: Turning Fear and Doubt into Fuel for Brilliance}
    
    ``\textit{Shatters misconceptions$\ldots$ Cain consistently holds the reader's interest by presenting individual profiles $\ldots$ and reporting on the latest studies. \textbf{Her diligence, research, and passion for this important topic has richly paid off.}}'' - Adam M. Grant, Ph.D., associate professor of management, the Wharton School of Business
\end{quotation}

\section*{Still more advance noise for \textit{quiet}}
\begin{quotation}
    ``\textit{\textbf{Once in a blue moon, a book comes along that gives us startling new insights}. \emph{Quiet} is that book: it's part page-turner, part cutting-edge science. The implications for business are especially valuable: \emph{Quiet} offers tips on how introverts can lead effectively, give winning speeches, avoid burnout, and choose the right roles. \textbf{This charming, gracefully written, thoroughly researched book is simply masterful.}}'' - Publishers Weekly
    
    ``\textit{\textbf{Quiet} elevates the conversation about introverts in our outwardly oriented society to new heights. \textbf{I think that many introverts will discover that, even though they didn't know it, they have been waiting for this book all their lives.}}'' - Adam S. McHugh, author of Introverts in the Church
    
    ``\textit{Susan Cain's \emph{Quiet} is wonderfully informative about the culture of the extrovert ideal and the psychology of a sensitive temperament, and she is helpfully perceptive about how introverts can make the most of their personality preferences in all aspects of life. \textbf{Society needs introverts, so everyone can benefit from the insights in this important book.}}'' - Jonathan M. Cheek, professor of psychology at Wellesley College, co-editor of \textit{Shyness: Perspectives on Research and Treatment}
    
    ``\textit{\textbf{A brilliant, important, and personally affecting book}. Cain shows that, for all its virtue, America's Extrovert Ideal takes up way too much oxygen. \textbf{Cain herself is the perfect person to make this case - with winning grace and clarity she shows us what it looks like to think outside the group.}}'' - Christine Kenneally, author of \textit{The First Word}
    
    ``\textit{What Susan Cain understands - and readers of this fascinating volume will soon appreciate - is something that psychology and our fast-moving and fast-talking society have been all too slow to realize: \textbf{Not only is there really nothing wrong with being quiet, reflective, shy, and introverted, but there are distinct advantages to being this way.}} - Jay Belsky, Robert M. and Natalie Reid Dorn Professor, Human and Community Development, University of California, Davis
    
    ``\textit{\textbf{Author Susan Cain exemplifies her own quiet power in this exquisitely written and highly reader page-turner.} She brings important research and poignant, personal examples into the light, greatly deepening our understanding of the introvert experience.''} - Jennifer B. Kahnweiler, Ph.D., author of \textit{The Introverted Leader}
    
    ``\textit{A species in which everyone was General Patton would not succeed, any more than would a race in which everyone was Vincent van Gogh. I prefer to think that \textbf{the planet needs athletes, philosophers, sex symbols,  painters, scientists; it needs the warmhearted, the hardhearted, the coldhearted, and the weakhearted}. It needs those who can devote their lives to studying how many droplets of water are secreted by the salivary glands of dogs under which circumstances, and it needs those who can  capture the passing impression of cherry blossoms in a fourteen-syllable poem or devote 25 pages to the dissection of a small boy's feelings as he lies in bed in the dark waiting for his mother to kiss him goodnight$\ldots$ Indeed the presence of outstanding strengths presupposes that energy needed in other areas has been channeled away from them.}'' - Allen Shawn
\end{quotation}

%------------------------------------------------------------------------------%

\section*{Author's Note}
\begin{quotation}
    ``I have been working on this book officially since 2005, and unofficially for my entire adult life.
    
    I have spoken and written to hundreds, perhaps thousands, of people about the topics covered inside, and have read as many books, scholarly papers, magazine articles, chat-room discussions, and blog posts.
    
    Some of these I mention in the book; others informed almost every sentence I wrote.
    
    \textit{Quiet} stands on many shoulders, especially the scholars and researchers whose work taught me so much.
    
    In a perfect world, I would have named every 1 of my sources, mentors, and interviewees.
    
    But for the sake of readability, some names appear
    only in the Notes or Acknowledgments.
    
    %
    For similar reasons, I did not use ellipses or brackets in certain quotations but made sure that the extra or missing words did not change the speaker's or writer's meaning.
    
    If you would like to quote these written sources from the original, the citations directing you to the full quotations appear in the Notes.
    
    %
    I've changed the names and identifying details of some of the people whose stories I tell, and in the stories of my own work as a lawyer and consultant.
    
    To protect the privacy of the participants in Charles di Cagno's public speaking workshop, who did not plan to be included in a book when they signed up for the class, the story of my 1st evening in class is a composite based on several sessions; so is the story of Greg and Emily, which is based on many interviews with similar couples.
    
    \fbox{Subject to the limitations of memory, all other stories are recounted as they happened or were told to me.}
    
    I did not fact-check the stories people told me about themselves, but only included those I believed to be true.''
\end{quotation}

%------------------------------------------------------------------------------%

\section*{\textsc{Introduction}: The North \& South of Temperament}

Montgomery, Alabama.

Dec 1, 1955.

Early evening.

A public bus pulls to a stop and a sensibly dressed woman in her 40s gets on.

She carries herself erectly, despite having spent the day bent over an ironing board in a dingy basement tailor shop at the Montgomery Fair department store.

Her feet are swollen, her shoulders ache.

She sits in the 1st row of the Colored section and watches quietly as the bus fills with riders.

Until the driver orders her to give her seat to a white passenger.

%
The woman utters a single word that ignites 1 of the most important civil rights protests of the 20th century, 1 word that helps America find its better self.

%
\fbox{The word is ``No.''}

%
The driver threatens to have her arrested.

%
``You may do that,'' says Rosa Parks.

%
A police officer arrives.

He asks Parks why she won't move.

%
``Why do you all push us around?'' she answers simply.

``I don't know,'' he says.

``But the law is the law, and you're under arrest.''

%
On the afternoon of her trial and conviction for disorderly conduct, the Montgomery Improvement Association holds a rally for Parks at the Holt Street Baptist Church, in the poorest section of town.

5000 gather to support Parks's lonely act of courage.

They squeeze inside the church until its pews can hold no more.

The rest wait patiently outside, listening through loudspeakers.

The Reverend Martin Luther King Jr. addresses the crowd.

``There comes a time that people get tired of being trampled over by the iron feet of oppression,'' he tells them.

``There comes a time when people get tired of being pushed out of the glittering sunlight of life's July and left standing amidst the piercing chill of an Alpine November.''

%
He praises Parks's bravery and hugs her.

She stands silently, her mere presence enough to galvanize the crowd.

The association launches a city-wide bus boycott that lasts 381 days.

The people trudge miles to work.

They carpool with strangers.

They change the course of American history.

%
Cain had always imagined Rosa Parks as a stately woman with a bold temperament, someone who could easily stand up to a busload of glowering passengers.

But when she died in 2005 at the age of 92, the flood of obituaries recalled her as soft-spoken, sweet, and small in stature.

They said she was \fbox{``timid and shy'' but had ``the courage of a lion.''}

They were full of phrases like \fbox{``radical humility'' and ``quiet fortitude.''}

What does it mean to be quiet \textit{and} have fortitude? these descriptions asked implicitly.

\fbox{How could you be shy \textit{and} courageous?}

%
Parks herself seemed aware of this paradox, calling her autobiography \textit{Quiet Strength} - a title that challenges us to question our assumptions.

Why \textit{shouldn't} quiet be strong?

And what else can quiet do that we don't give it credit for?
\begin{center}
    $\star$
\end{center}
\fbox{Our lives are shaped as profoundly by personality as by gender or race.}

And the single most important aspect of personality - the ``north and south of temperament,'' as 1 scientist puts it - is where we fall on the \fbox{introvert-extrovert spectrum}.

Our place on this continuum influences our choice of friends and mates, and how we make conversation, resole differences, and show love.

It affects the careers we choose and whether or not we succeed at them.

It governs how likely we are to exercise, commit adultery, function well without sleep, learn from our mistakes, place big bets in the stock market, delay gratification, be a good leader, and ask ``what if.''\footnote{\textbf{Answer key.} exercise: extroverts; commit adultery: extroverts; function well without sleep: introverts; learn from our mistakes: introverts; place big bets: extroverts; delay gratification: introverts; be a good leader: in some cases introverts, in other cases extroverts, depending on the type of leadership called for; ask ``what if'': introverts.}

It's reflected in our brain pathways, neurotransmitters, and remote corners of our nervous systems.

Today introversion and extroversion are 2 of the most exhaustively researched subjects in \fbox{personality psychology}, arousing the curiosity of hundreds of scientists.

%
These researchers have been exciting discoveries aided by the latest technology, but they're part of a long and storied tradition.

Poets and philosophers have been thinking about introverts and extroverts since the dawn of recorded time.

Both personality types appear in the Bible and in the writings of Greek and Roman physicians, and some evolutionary psychologists say that the history of these types reaches back even farther than that: the animal kingdom also boasts ``introverts'' and ``extroverts,'' and we'll see, from fruit flies to pumpkinseed fish to rhesus monkeys.

As with other complementary pairings - masculinity and femininity, East and West, liberal and conservative - humanity would be unrecognizable, and vastly diminished, without both personality styles.

%
Take the partnership of Rosa Parks and Martin Luther King Jr.: a formidable orator refusing to give up his seat on a segregated bus wouldn't have had the same effect as a modest woman who'd clearly prefer to keep silent but for the exigencies of the situation.

And Parks didn't have the stuff to thrill a crowd if she'd tried to stand up and announce that she had a dream.

But with the King's help, she didn't have to.

%
Yet today we make room for a remarkably narrow range of personality styles.

We're told that to be great is to be bold, to be happy is to be sociable.

We see ourselves as a nation of extroverts - which means that we've lost sight of who we really are.

Depending on which study you consult, 1 3rd to 1 half of Americans are introverts - in other words, \textit{1 out of every 2 or 3 people you know}.

(Given that the United States is among the most extroverted of nations, the number must be at least as high in other parts of the world.)

If you're not an introvert yourself, you are surely raising, managing, married to, or coupled with one.

%
If these statistics surprise you, that's \fbox{probably because so many people pretend to be extroverts}.

Closet introverts pass undetected on playgrounds, in high school locker rooms, and in the corridors of corporate America.

Some fool even themselves, until some life event - a layoff, an empty nest, an inheritance that frees them to spend time as they like - jolts them into taking stock of their true natures.

You have only to raise the subject of this book with your friends and acquaintances to find that the most unlikely people consider themselves introverts.

%
\fbox{It makes sense that so many introverts hide even from themselves.}

We live with a value system that Cain calls the Extrovert Ideal - the omnipresent belief that the ideal self is gregarious, alpha, and comfortable in the spotlight.

\fbox{The archetypal extrovert prefers action to contemplation, risk-taking to heed-taking, certainty to doubt.}

He favors quick decisions, even at the risk of being wrong.

She works well in teams and socializes in groups.

We like to think that we value individuality, but all too often we admire 1 \textit{type} of individual - the kind who's comfortable ``putting himself out there.''

Sure, we allow technologically gifted loners who launch companies in garages to have any personality they please, but they are the exceptions, not the rule, and our tolerance extends mainly to those who get fabulously wealthy or hold the promise of doing so.

%
Introversion - along with its cousins sensitivity, seriousness, and shyness - is now a 2nd-class personality trait, somewhere between a disappointment and a pathology.

Introverts living under the Extrovert Ideal are like women in a man's world, discounted because of a trait that goes to the core of who they are.

Extroversion is an enormously appealing personality style, but we've turned it into an oppressive standard to which most of us feel we must conform.

%
The Extrovert Ideal has been documented in many studies, though this research has never been grouped under a single name.

Talkative people, e.g., are rated as smarter, better-looking, more interesting, and more desirable as friends.

Velocity of speech counts as well as volume: we rank fast talkers as more competent and likable than slow ones.

The same dynamics apply in groups, where research shows that the voluble are considered smarter than the reticent - even though there's zero correlation between the gift of gab and good ideas.

Even the word \textit{introvert} is stigmatized - 1 informal study, by psychologist Laurie Helgoe, found that introverts described their own physical appearance in vivid language (``green-blue eyes,'' ``exotic,'' ``high cheekbones''), but when asked to describe generic introverts they drew a bland and distasteful picture (``ungainly,'' ``neutral colors,'' ``skin problems'').

%
But we make a grave mistake to embrace the Extrovert Ideal so unthinkingly.

Some of our greatest ideas, art, and inventions - from the theory of evolution to van Gogh's sunflowers to the personal computer - came from quiet and cerebral people who knew how to tune in to their inner worlds and the treasures to be found there.

Without introverts, the world would be devoid of:
\begin{itemize}
    \item the theory of gravity
    \item the theory of relativity
    \item W. B. Yeats's ``The Second Coming''
    \item Chopin's nocturnes
    \item Proust's \textit{In Search of Lost Time}
    \item Peter Pan
    \item Orwell's \textit{1984} and \textit{Animal Farm}
    \item The Cat in the Hat
    \item Charlie Brown
    \item \textit{Schinder's List, E. T., and \textit{Close Encounters of the 3rd Kind}}
    \item Google
    \item Harry Potter\footnote{Sir Isaac Newton, Albert Einstein, W. B. Yeats, Frédéric Chopin, Marcel Proust, J. M. Barrie, George Orwell, Theodor Geisel (Dr. Seuss), Charles Schulz, Steven
Spielberg, Larry Page, J. K. Rowling.}
\end{itemize}
%
As the science journalist Winifred Gallagher writes: ``The glory of the disposition that stops to consider stimuli rather than rushing to engage with them is its \fbox{long association with intellectual and artistic achievement}.

Neither $E = mc^2$ nor \textit{Paradise Lost} was dashed off by a party animal.''

Even in less obviously introverted occupations, like finance, politics, and activism, some of the greatest leaps forward were made by introverts.

In this book we'll see how figures like Eleanor Roosevelt, Al Gore, Warren Buffett, Gandhi - and Rosa Parks - achieved what they did \fbox{not in spite of but \textit{because} of} their introversion.

%
Yet, as \textit{Quiet} will explore, many of the most important institutions of contemporary life are designed for those who enjoy group projects and high levels of stimulation.

As children, our classroom desks are increasingly arranged in pods, the better to foster group leaning, and research suggests that the vast majority of teachers believe that the ideal student is an extrovert.

We watch TV shows whose protagonists are not the ``children next door,'' like the Cindy Bradys and Beaver Cleavers of yesteryear, but rock stars and webcast hostesses with outsized personalities, like Hannah Montana and Carly Shay of \textit{iCarly}.

Even Sid the Science Kid, a PBS-sponsored role model for the preschool set, kicks off each school day by performing dance moves with his pals.

(``Check out my moves! I'm a rock star!'')

%
As adults, many of us work for organizations that insist we work in teams, in offices without walls, for supervisors who value ``people skills'' above all.

To advance our careers, we're expected to promote ourselves unabashedly.

The scientists whose research gets funded often have confident, perhaps overconfident, personalities.

The artists whose work adorns the walls of contemporary museums strike impressive poses at gallery openings.

The authors whose books get published - once accepted as a reclusive breed - are now vetted by publicists to make sure they're talk-show ready.

(You wouldn't be reading this book if Cain hadn't convinced her publisher that she was enough of a pseudo-extrovert to promote it.)

%
If you're an introvert, you also know that \fbox{the bias against quiet can cause deep psychic pain.}

As a child you might have overheard your parents apologize for your shyness.

(``Why can't you be more like the Kennedy boys?'' the Camelot-besotted parents of 1 man Cain interviewed repeatedly asked him.)

Or at school you might have been prodded to come ``out of your shell'' - that noxious expression which fails to appreciate that some animals naturally carry shelter everywhere they go, and that some humans are just the same.

``All the comments from childhood still ring in my ears, that I was lazy, stupid, slow, boring,'' writes a member of an e-mail list called Introvert Retreat.

``By the time I was old enough to figure out that I was simply introverted, it was a part of my being, the assumption that there is something inherently wrong with me.

I wish I could find that little vestige of doubt and remove it.''

%
Now that you're an adult, you might still feel a pang of guilt when you decline a dinner invitation in favor of a good book.

Or maybe you like to eat alone in restaurants and could do without the pitying looks from fellow diners.

Or you're told that you're ``\fbox{in your head too much},'' a phrase that's often deployed against the quiet and cerebral.

%
Of course, there's another word for such people: \fbox{thinkers}.
\begin{center}
    $\star$
\end{center}
Cain has seen firsthand \textit{how difficult it is for introverts to take stock of their own talents, and how powerful it is when finally they do}.

For more than 10 years Cain trained people of all stripes - corporate lawyers and college students, hedge-fund managers and married couples - in negotiation skills.

Of course, we covered the basics: how to prepare for a negotiation, when to make the 1st offer, and what to do when the other person says ``take it or leave it.''

But Cain also helped clients figure out their natural personalities and how to make the most of them.

%
Cain's very 1st client was a young woman named Laura.

She was a Wall Street lawyer, but a quiet and daydreamy one who dreaded the spotlight and disliked aggression.

She had managed somehow to make it through the crucible of Harvard Law School - a place where classes are conducted in huge, gladiatorial amphitheaters, and where she once got so nervous that she threw up on the way to class.

Now that she was in the real world, she wasn't sure she could represent her clients as forcefully as they expected.

%
For the 1st 3 years on the job, Laura was so junior that she never had to test this premise.

But 1 day the senior lawyer she's been working with went on vacation, leaving her in charge of an important negotiation.

The client was a South American manufacturing company that was about to default on a bank loan and hoped to renegotiate its terms; a syndicate of bankers that owned the endangered loan sat on the other side of the negotiating table.

%
Laura would have preferred to hide under said table, but she was accustomed to fighting such impulses.

Gamely but nervously, she took her spot in the dead chair, flanked by her clients: general counsel on 1 side and senior financial officer on the other.

These happened to be Laura's favorite clients: gracious and soft-spoken, very different from the master-of-the-universe types her firm usually represented.

In the past, Laura had taken the general counsel to a Yankees game and the financial officer shopping for a handbag for her sister.

But now these cozy outings - just the kind of socializing Laura enjoyed - seemed a world away.

Across the table sat 9 disgruntled investment bankers in tailored suits and expensive shoes, accompanied by their lawyer, a square-jawed woman with a hearty manner.

Clearly not the self-doubting type, this woman launched into an impressive speech on how Laura's clients would be lucky simply to accept the bankers' terms.

It was, she said, a very magnanimous offer.

%
Everyone waited for Laura to reply, but she couldn't think of anything to say.

So she just sat there.

Blinking.

All eyes on her.

Her clients shifting uneasily in their seats.

Her thoughts running in a familiar loop: \textit{I'm too quiet for this kind of thing, too unassuming, too cerebral}.

She imagined the person who would be better equipped to save the day: someone bold, smooth, ready to pound the table.

In middle school this person, unlike Laura, would have been called ``outgoing,'' the highest accolade her 7th-grade classmates knew, higher even than ``pretty,'' for a girl, or ``athletic,'' for a guy.

Laura promised herself that she only had to make it through the day.

Tomorrow she would go look for another career.

%
Then she remembered what Cain's told her again and again: she was an introvert, and as such she had unique powers in negotiation - perhaps less obvious but no less formidable.

She'd probably prepared more than everyone else.

She had a \fbox{quiet but firm speaking style}.

She \fbox{rarely spoke without thinking}.

Being mild-mannered, she could take strong, even aggressive, positions while coming across as perfectly reasonable.

And she tended to ask questions - lots of them - and actually listen to the answers, which, no matter what your personality, is crucial to strong negotiation.

%
So Laura finally stared doing what came naturally.

%
``Let's go back a step.

What are your numbers based on?'' she asked.

%
``What if we structured the loan this way, do you think it might work?''

%
``That way?''

%
``Some other way?''

%
At 1st her questions were tentative.

She picked up steam as she went along, posing them more forcefully and making it clear that she's done her homework and wouldn't concede the facts.

But she also stayed true to her own style, never raising her voice or losing her decorum.

Every time the bankers made an assertion that seemed unbudgeable, Laura tried to be constructive.

``Are you saying that's the only way to go?

What if we took a different approach?"

%
Eventually her simple queries shifted the mood in the room, just as the negotiation textbooks say they will.

The bankers stopped speechifying and dominance-posing, activities for which Laura felt hopelessly ill-equipped, and they started having an actual conversation.

%
More discussion.

Still no agreement.

1 of the bankers revved up again, throwing his papers down and storming out of the room.

Laura ignored this display, mostly because she didn't know what else to do.

Later on someone told her that at that pivotal moment she'd played a good game of something called ``negotiation jujitsu''; but she knew that she was just doing what you learn to do naturally as a quiet person in a loudmouth world.

%
Finally the 2 sides struck a deal.

The bankers left the building, Laura's favorite clients headed for the airport, and Laura went home, curled up with a book, and tried to forget the day's tensions.

%
But the next morning, the lead lawyer for the bankers - the vigorous woman with the strong jaw - called to offer her a job.

``I've never seen anyone so nice and so tough at the same time,'' she said.

And the day after that, the lead banker called Laura, asking if \textit{her} law firm would represent \textit{his} company in the future.

``We need someone who can help us put deals together without letting ego get in the way,'' he said.

%
By sticking to her own gentle way of doing things, Laura had reeled in new business for her firm and a job offer for herself.

Raising her voice and pounding the table was unnecessary.

%
Today Laura understands that her introversion is an essential part of who she is, and she embraces her reflective nature.

The loop inside her head that accused her of being too quiet and unassuming plays much less often.

Laura knows that she can hold her own when she needs to.
\begin{center}
    $\star$
\end{center}
``What exactly do I mean when I say that Laura is an \textit{introvert}?'' - Cain wonders.

When Cain started writing this book, the 1st thing Cain wanted to find out was precisely how researchers define introversion and extroversion.

Cain knew that in 1921 the influential psychologist Carl Jung had published a bombshell of a book, \textit{Psychological Types}, popularizing the terms \textit{introvert} and \textit{extrovert} as the central building blocks of personality.

\textit{Introverts are drawn to the inner world of thought and feeling}, said Jung, \textit{extroverts to the external life of people and activities}.

Introverts focus on the meaning they make of the events swirling around them; extroverts plunge into the events themselves.

Introverts recharge their batteries by being alone; extroverts need to recharge when they don't socialize enough.

If you've ever taken a Myers-Briggs personality test, which is based on Jung's thinking and used by the majority of universities and Fortune 100 companies, then you may already be familiar with these ideas.

%
But what do contemporary researchers have to say?

Cain soon discovered that there is no all-purpose definition of introversion or extroversion; these are not unitary categories, like ``curly-haired'' or ``16-year-old,'' in which everyone can agree on who qualifies for inclusion.

E.g., adherents of the Big 5 school of personality psychology (which argues that human personality can be boiled down to 5 primary traits) define introversion not in terms of a rich inner life but as a lack of qualities e.g. assertiveness and sociability.

There are almost as many definitions of \textit{introvert} and \textit{extrovert} as there are personality psychologists, who spend a great deal of time arguing over which meaning is most accurate.

Some think that Jung's ideas are outdated; others swear that he's the only one who got it right.

%
Still, today's psychologists tend to agree on several important point: e.g., that introverts and extroverts differ in the level of outside stimulation that they need to function well.

Introverts feel ``just right'' with less stimulation, as when they sip wine with a close friend, solve a crossword puzzle, or read a book.

Extroverts enjoy the extra bang that comes from activities like meeting new people, skiing slippery slopes, and cranking up the stereo.

``Other people are very arousing,'' says the personality psychologist David Winter, explaining why your typical introvert would rather spend her vacation reading on the beach than partying on a cruise ship.

``They arouse threat, fear, flight, and love.

100 people are very stimulating compared to 100 books or 100 grains of sand.''

%
Many psychologists would also agree that introverts and extroverts work differently.

Extroverts tend to tackle assignments quickly.

They make fast (sometimes rash) decisions, and are comfortable multitasking and risk-taking.

They enjoy ``the thrill of the chase'' for rewards like money and status.

%
\fbox{Introverts often work more slowly and deliberately.}

They like to \fbox{focus on 1 task at a time and can have mighty powers of concentration.}

They're \fbox{relatively immune to the lures of wealth and fame.}

%
\fbox{Our personalities also shape our social styles.}

Extroverts are the people who will add life to your dinner party and laugh generously at your jokes.

They tend to be assertive, dominant, and in great need of company.

Extroverts think out loud and on their feet; they prefer talking to listening, rarely find themselves at a loss for words, and occasionally blurt out things they never meant to say.

They're \fbox{comfortable with conflict, but not with solitude.}

%
Introverts, in contrast, may have strong social skills and enjoy parties and business meetings, but after a while wish they were home in their pajamas.

They prefer to devote their social energies to close friends, colleagues, and family.

They \textit{listen more than they talk, think before they speak}, and often feel as if they \textit{express themselves better in writing than in conversation}.

They tend to dislike conflict.

Many have a horror of small talk, but enjoy deep discussions.

%
A few things introverts are not: The word \textit{introvert} is not a synonym for hermit or misanthrope.

Introverts \textit{can} be these things, but most are perfectly friendly.

1 of the most humane phrases in the English language - ``Only connect!'' - was written by the distinctly introverted E. M. Forster in a novel exploring the question of how to achieve ``human love at its height.''

%
\fbox{Nor are introverts necessarily shy.}

Shyness is the fear of social disapproval or humiliation, while introversion is a preference for environments that are not overstimulating.

Shyness is inherently painful; introversion is not.

1 reason that people confuse the 2 concepts is that they sometimes overlap (though psychologists debate to what degree).

Some psychologists map the 2 tendencies on vertical and horizontal axes, with the introvert-extrovert spectrum on the horizontal axis, and the anxious-stable spectrum on the vertical.

With this model, you end up with 4 quadrants of personality types: \textit{calm extroverts, anxious} (or \textit{impulsive}) \textit{extroverts, calm introverts}, and \textit{anxious introverts}.

In other words, you can be a shy extrovert, like Barbra Streisand, who has a larger-than-life personality and paralyzing stage fright; or a non-shy introvert, like Bill Gates, who by all accounts keeps to himself but is unfazed by the opinions of others.

%
You can also, of course, be both shy \textit{and} an introvert: T. S. Eliot was a famously private soul who wrote in ``The Waste Land'' that he could ``show you fear in a handful of dust.''

Many shy people turn inward, partly as a refuge from the socializing that causes them such anxiety.

And many introverts are shy, partly as a result of receiving the message that there's something wrong with their preference for reflection, and partly because their physiologies, as we'll see, compel them to withdraw from high-stimulation environments.

%
But for all their differences, shyness and introversion have in common something profound.

The mental state of a shy extrovert sitting quietly in a business meeting may be very different from that of a calm introvert - the shy person is afraid to speak up, while the introvert is simply overstimulated - but to the outside world, the 2 appear to be the same.

This can give both types insight into how our reverence for alpha status blinds us to things that are good and smart and wise.

For very different reasons, shy and introverted people might choose to spend their days in behind-the-scenes pursuits like inventing, or researching, or holding the hands of the gravely ill - or in leadership positions they execute with quiet competence.

These are not alpha roles, but the people who play them are role models all the same.
\begin{center}
    $\star$
\end{center}
If you're still not sure where you fall on the introvert-extrovert spectrum, you can access yourself here.

Answer each question ``true'' or ``false,'' choosing the answer that applies to you more often than not.\footnote{This is an informal quiz, not a scientifically validated personality test.

The questions were formulated based on characteristics of introversion often accepted by contemporary researchers.}
\begin{enumerate}
    \item I prefer 1-on-1 conversations to group activities.
    \item I often prefer to express myself in writing.
    \item I enjoy solitude.
    \item I seem to care less than my peers about wealth, fame, and status.
    \item I dislike small talk, but I enjoy talking in depth about topics that matter to me.
    \item People tell me that I'm a good listener.
    \item I'm not a big risk-taker.
    \item I enjoy work that allows me to ``dive in'' with few interruptions.
    \item I like to celebrate birthdays on a small scale, with only 1 or 2 close friends or family members.
    \item People describe me as ``soft-spoken'' or ``mellow''.
    \item I prefer not to show or discuss my work with others until it's finished.
    \item I dislike conflict.
    \item I do my best work on my own.
    \item I tend to think before I speak.
    \item I feel drained after being out and about, even if I've enjoyed myself.
    \item I often let calls go through to voice mail.
    \item If I had to choose, I'd prefer a weekend with absolutely nothing to do to one with too many things scheduled.
    \item I don't enjoy multitasking.
    \item I can concentrate easily.
    \item In classroom situations, I prefer lectures to seminars.
\end{enumerate}
%
The more often you answered ``true,'' the more introverted you probably are.

If you found yourself with a roughly equal number of ``true'' and ``false'' answers, then you may be an ambivert - yes, there really is such a word.

%
But even if you answered every single question as an introvert or extrovert, that doesn't mean that your behavior is predictable across all circumstances.

We can't say that every introvert is a bookworm or every extrovert wears lampshades at parties any more than we can say that every woman is a natural consensus-builder and every man loves contact sports.

As Jung felicitously put it, ``There is no such thing as a pure extrovert or a pure introvert.

Such a man would be in the lunatic asylum.''

%
This is partly because we are all \fbox{gloriously complex individuals}, but also because there are so many different \textit{kinds} of introverts and extroverts.

Introversion and extroversion interact with our other personality traits and personal histories, producing wildly different kinds of people.

So if you're an artistic American guy whose father wished you'd try out for the football team like your rough-and-tumble brothers, you'll be a very different kind of introvert from, say, a Finnish businesswoman whose parents were lighthouse keepers.

(Finland is a famously introverted nation.

Finnish joke: How can you fell if a Finn likes you?

He's staring at your shoes instead of his own.)

%
\fbox{Many introverts are also ``highly sensitive,''} which sounds poetic, but is actually a technical term in psychology.

If you are a sensitive sort, then you're more apt than the average person to feel pleasantly overwhelmed by Beethoven's ``Moonlight Sonata'' or a well-turned phrase or an act of extraordinary kindness.

You may be quicker than others to feel sickened by violence and ugliness, and you likely have a very strong conscience.

When you were a child you were probably called ``shy,'' and to this day feel nervous when you're being evaluated, e.g. when giving a speech or on a 1st date.

Later we'll examine why this seemingly unrelated collection of attributes tends to belong to the same person and why this person is often introverted.

(No one knows exactly how many introverts are highly sensitive, but we know that 70\% of sensitives are introverts, and the other 30\% tend to report needing a lot of ``down time.'')

%
All of this complexity means that not everything you read in \textit{Quiet} will apply to you, even if you consider yourself a true-blue introvert.

For 1 thing, we'll spend some time talking about shyness and sensitivity, while you might have neither of these traits.

That's OK.

Take what applies to you, and use the rest to improve your relationships with others.

%
Having said all this, in \textit{Quiet} we'll try not to get too hung up on definitions.

Strictly defining terms is vital for researchers whose studies depend on pinpointing exactly where introversion stops and other traits, like shyness, start.

But in \textit{Quiet} we'll concern ourselves more with the \textit{fruit} of that research.

Today's psychologists, joined by neuroscientists with their brain-scanning machines, have unearthed illuminating insights that are changing the way we see the world - and ourselves.

They are answering questions e.g.: \textit{Why are some people talkative while others measure their words?}

\textit{Why do some people burrow into their work and others organize office birthday parties?}

\textit{Why are some people comfortable wielding authority while others prefer neither to lead nor to be led?}

\textit{Can} introverts be leaders?

Is our cultural preference for extroversion in the natural order of things, or is it socially determined?

\fbox{From an evolutionary perspective, introversion must have survived as a personality trait for a reason} - so what might the reason be?

If you're an introvert, should you devote your energies to activities that come naturally, or should you stretch yourself, as Laura did that day at the negotiation table?

%
The answers might surprise you.

%
If there is only 1 insight you take away from this book, though, Cain hopes it's a newfound sense of entitlement to be yourself.

Cain can vouch personally for the life-transforming effects of this outlook.

Remember that 1st client Cain told you about, the one Cain called Laura in order to protect her identity?

%
That was a story about Cain herself.

Cain was her own 1st client.

%------------------------------------------------------------------------------%

\part{\textsc{The Extrovert Ideal}}

\section{\textsc{The Rise of the ``Mighty Likable Fellow''}: How Extroversion Became the Cultural Ideal}

\begin{quotation}
    \textit{Strangers' eyes, keen and critical}.
    
    \textit{Can you met them proudly - confidently - without fear?} - Print Advertisement for Woodbury's soap, 1922
\end{quotation}
The date: 1902.

The place: Harmony Church, Missouri, a tiny, dot-on-the-map town located on a floodplain 100 miles from Kansas City.

Our young protagonist: a good-natured but insecure high school student named Dale.

%
Skinny, unathletic, and fretful, Dale is the son of a morally upright but perpetually bankrupt pig farmer.

He respects his parents but dreads following in their poverty-stricken footsteps.

Dale worries about other things, too: thunder and lighting, going to hell, and being tongue-tied at crucial moments.

He even fears his wedding day: What if he can't think of anything to say to his future bride?

%
1 day a Chautauqua speaker comes to town.

The Chautauqua movement, born in 1873 and based in upstate New York, sends gifted speakers across the country to lecture on literature, science, and religion.

Rural Americans prize these presenters for the whiff of glamour they bring from the outside world - and their power to mesmerize and audience.

This particular speaker captivates the young Dale with his own rags-to-riches tale: once he'd been a lowly farm boy with a bleak future, but he developed a charismatic speaking style and took the stage at Chautauqua.

Dale hangs on his every word.

%
A few years later, Dale is again impressed by the value of public speaking.

His family moves to a farm 3 miles outside of Warrensburg, Missouri, so he can attend college there without paying room and board.

Dale observes that the students who win campus speaking contests are seen as leaders, and he resolves to be 1 of them.

he signs up for every contest and rushes home at night to practice.

Again and again he loses; Dale is dogged, but not much of an orator.

Eventually, though, his efforts begin to pay off.

He transforms himself into a speaking champion and campus hero.

Other students turn to him for speech lessons; he trains them and they start winning, too.

%
By the time Dale leaves college in 1908, his parents are still poor, but corporate America is booming.

Henry Ford is selling Model Ts like griddle cakes, using the slogan ``for business and for pleasure.''

J. C. Penney, Woolworth, and Sears Roebuck have become household names.

Electricity lights up the homes of the middle class; indoor plumbing spares them midnight trips to the outhouse.

%
The new economy calls for a new kind of man - a salesman, a social operator, someone with a ready smile, a masterful handshake, and the ability to get along with colleagues while simultaneously outshining them.

Dale joins the swelling ranks of salesmen, heading out on the road with few possessions but his \fbox{silver tongue}.

%
Dale's last name is Carnegie (Carnagey, actually; he changes the spelling later, likely to evoke Andrew, the great industrialist).

After a few grueling years selling beef for Armour and Company, he sets up shop as a public-speaking teacher.

Carnegie holds his 1st class at a YMCA night school on 125th Street in New York City.

He asks for the usual 2-dollars-per-session salary for night school teachers.

The Y's director, doubting that a public-speaking class will generate much interest, refuses to pay that kind of money.

%
But the class is an overnight sensation, and Carnegie goes on to found the Dale Carnegie Institute, dedicated to helping businessmen root out the very insecurities that had held him back as a young man.

In 1913 he publishes his 1st book, \textit{Public Speaking and Influencing Men in Business}.

``In the days when pianos and bathrooms were luxuries,'' Carnegie writes, ``men regarded ability in speaking as a peculiar gift, needed only by the lawyer, clergy man, or statesman.

Today we have come to realize that it is the indispensable weapon of those who would forge ahead in the keen competition of business.''
\begin{center}
    $\star$
\end{center}
Carnegie's metamorphosis from farmboy to salesman to public-speaking icon is also the story of the rise of the Extrovert Ideal.

Carnegie's journey reflected a cultural evolution that reached a tipping point around the turn of the 20th century, changing forever who we are and whom we admire, how we act at job interviews and what we look for in an employee, how we court our mates and raise our children.

America had shifted from what the influential cultural historian Warren Susman called a Culture of Character to a Culture of Personality - and opened up a Pandora's Box of personal anxieties from which we would never quite recover.

%
\fbox{In the Culture of Character, the ideal self was serious, disciplined, and honorable.}

What counted was not so much the impression one made in public as how one behaved in private.

The word \textit{personality} didn't exist in English until the 18th century, and the idea of ``having a good personality'' was not widespread until the 20th.

%
But when they embraced the Culture of Personality, Americans started to focus on how others perceived them.

They became captivated by people who were bold and entertaining.

``The social role demanded of all in the new Culture of Personality was that of a performer,'' Susman famously wrote.

``Every American was to become a performing self.''

%
The rise of industrial America was a major force behind this cultural evolution.

The nation quickly developed from an agricultural society of little houses on the prairie to an urbanized, ``the business of America is business'' powerhouse.

In the country's early days, most Americans lived like Dale Carnegie's family, on farms or in small towns, interacting with people they'd known since childhood.

But when the 20th century arrived, a perfect storm of big business, urbanization, and mass immigration blew the population into the cities.

In 1790, only 3\% of Americans lived in cities; in 1840, only 8\% did; by 1920, more than a 3rd of the country were urbanites.

``We cannot all live in cities,'' wrote the news editor Horace Greeley in 1867, ``yet nearly all seem determined to do so.''

%
Americans found themselves working no longer with neighbors but with strangers.

``Citizens'' morphed into ``employees,'' facing the question of how to make a good impression on people to whom they had no civic or family ties.

``The reasons why 1 man gained a promotion or 1 woman suffered a social snub,'' writes the historian Roland Marchand, ``had become less explicable on grounds of long-standing favoritism or old family feuds.

In the increasingly anonymous business and social relationships of the age, one might suspect that anything - including a 1st impression - had made the crucial difference.''

Americans responded to these pressures by trying to become salesmen who could sell not only their company's latest gizmo but also themselves.

%
1 of the most powerful lenses through which to view the transformation from Character to Personality is the self-help tradition in which Dale Carnegie played such a prominent role.

Self-help books have always loomed large in the American psyche.

Many of the earliest conduct guides were religious parables, like \textit{The Pilgrim's Progress}, published in 1678, which warned readers to behave with restraint if they wanted to make it into heaven.

The advice manuals of the 19th century were less religious but still preached the value of a noble character.

They featured case studies of historical heroes like Abraham Lincoln, revered not only as a gifted communicator but also as a modest man who did not, as Ralph Waldo Emerson put it, ``offend by superiority.''

They also celebrated regular people who lived highly moral lives.

A popular 1899 manual called \textit{Character: The Grandest Thing in the World} featured a timid shop girl who gave away her meager earnings to a freezing beggar, then rushed off before anyone could see what she's done.

Her virtue, the reader understood, derived not only from her generosity but also from her wish to remain anonymous.

%
But by 1920, popular self-help guides had changes their focus from inner virtue to outer charm - ``to know \textit{what} to say and \textit{how} to say it,'' as 1 manual put it.

``To create a personality is power,'' advised another.

``Try in every way to have a ready command of the manners which make people think `he's a mighty likable fellow,''' said a 3rd.

``That is the beginning of a reputation for personality.''

\textit{Success} magazine and \textit{The Saturday Evening Post} introduced departments instructing readers on the art of conversation.

The same author, Orison Swett Marden, who wrote \textit{Character: The Grandest Thing in the World} in 1899, produced another popular title in 1921.

It was called \textit{Masterful Personality}.

%
Many of these guides were written for businessmen, but women were also urged to work on a mysterious quality called ``fascination.''

Coming of age in the 1920s was such a competitive business compared to what their grandmothers had experienced, warned 1 beauty guide, that they had to be visibly charismatic: ``People who pass us on the street can't know that we're clever and charming unless we look it.''

%
Such advice - ostensibly meant to improve people's lives - must have made even reasonably confident people uneasy.

Susman counted the words that appeared most frequently in the personality-driven advice manuals of the early 20th century and compared them to the character guides of the 19th century.

The earlier guides emphasized attributes that anyone could work on improving, described by words like
\begin{example}
    Citizenship, Duty, Work, Golden deeds, Honor, Reputation, Morals, Manners, Integrity
\end{example}
%
But the new guides celebrated qualities that were - no matter how easy Dale Carnegie made it sound - trickier to acquire.

Either you embodied these qualities or you didn't:
\begin{example}
    Magnetic, Fascinating, Stunning, Attractive, Glowing, Dominant, Forceful, Energetic
\end{example}
%
It was no coincidence that in the 1920s and the 1930s, American became obsessed with movie stars.

Who better than a matinee idol to model personal magnetism?
\begin{center}
    $\star$
\end{center}
Americans also received advice on self-presentation - whether they liked it or not - from the advertising industry.

While early print ads were straightforward product announcements (``\textsc{Eaton's highland linen: the freshest and cleanest writing paper}''), the new personality-driven ads cast consumers as performers with stage fright from which only the advertiser's product might rescue them.

These ads focused obsessively on the hostile glare of the public spotlight.

``\textsc{all around you people are judging you silently},'' warned a 1922 ad for Woodbury's soap.

``\textsc{critical eyes are sizing you up right now},'' advised the Williams Shaving Cream company.

%
Madison Avenue spoke directly to the anxieties of male salesmen and middle managers.

In 1 ad for Dr. West's toothbrushes, a prosperous-looking fellow sat behind a desk, his arm cocked confidently behind his hip, asking whether you've ``\textsc{ever tried selling yourself to you? a favorable 1st impression is the greatest single factor in business or social success}.''

The Williams Shaving Cream ad featured a slick-haired, mustachioed man urging readers to ``\textsc{let your face reflect confidence, not worry! it's the `look' of you by which you are judged most often}.''

%
Other ads reminded women that their success in the dating game depended not only on looks but also on personality.

In 1921 a Woodbury's soap ad showed a crestfallen young woman, home alone after a disappointing evening out.

She had ``longed to be successful, gay, triumphant,'' the text sympathized.

But without the help of the right soap, the woman was a social failure.

%
10 years later, Lux laundry detergent ran a print ad featuring a plaintive letter written to Dorothy Dix, the Dear Abby of her day.

``Dear Miss Dix,'' read the letter, ``How can I make myself more popular?

I am fairly pretty and not a dumbbell, but I am so timid and self-conscious with people.

I'm always sure they're not going to like me$\ldots$. - Joan G.''

%
Miss Dix's answer came back clear and firm.

If only Joan would use Lux detergent on her lingerie, curtains, and sofa cushions, she would soon gain a ``deep, sure, inner conviction of being charming.''

%
This portrayal of courtship as a high-stakes performance reflected the bold new mores of the Culture of Personality.

Under the restrictive (in some cases repressive) social codes of the Culture of Character, both genders displayed some reserve when it came to the mating dance.

Women who were too cloud or made inappropriate eye contact with strangers were considered brazen.

Upper-class women had more license to speak than did their lower-class counterparts, and indeed were judged partly on their talent for witty repartee, but even they were advised to display blushes and downcast eyes.

They were warned by conduct manuals that ``the coldest reserve'' was ``more admirable in a woman a man wishe[d] to make his wife than the least approach to undue familiarity.''

Men could adopt a quiet demeanor that implied self-possession and a power that didn't need to flaunt itself.

Though shyness per se was unacceptable, reserve was a mark of good breeding.

%
But with the advent of the Culture of Personality, the value of formality began to crumble, for women and men alike.

Instead of paying ceremonial calls on woman and making serious declarations of intentions, men were now expected to launch verbally sophisticated courtships in which they threw women ``a line'' of elaborate flirtatiousness.

Men who were too quiet around women risked being thought gay; as a popular 1926 sex guide observed, ``homosexuals are invariably timid, shy, retiring.''

Women, too, were expected to walk a fine line between propriety and boldness.

If they responded to shyly to romantic overtures, they were sometimes called ``frigid.''

%
The field of psychology also began to grapple with the pressure to project confidence.

In the 1920s an influential psychologist named Gordon Allport created a diagnostic test of ``Ascendance-Submission'' to measure social dominance.

``Our current civilization,'' observed Allport, who was himself shy and reserved, ``seems to place a premium upon the aggressive person, the `go-getter.'''

In 1921, Carl Jung noted the newly precarious status of introversion.

Jung himself saw introverts as ``educators and promoters of culture'' who showed the value of ``the interior life which is so painfully wanting in our civilization.''

But he acknowledged that their ``reserve and apparently groundless embarrassment naturally arouse all the current prejudices against this type.''

%
But nowhere was the need to appear self-assured more apparent than in a new concept in psychology called the \textit{inferiority complex}.

The IC, as it became known in the popular press, was developed in the 1920s by a Viennese psychologist named Alfred Adler to describe feelings of inadequacy and their consequences.

``Do you feel insecure?'' inquired the cover of Adler's best-selling book, \textit{Understanding Human Nature}.

``Are you fainthearted?

Are you submissive?''

Adler explained that all infants and small children feel inferior, living as they do in a world of adults and older siblings.

In the normal process of growing up they learn to direct these feelings into pursuing their goals.

But if things go awry as they mature, they might be saddled with the dreaded IC - a grave liability in an increasingly competitive society.

%
The idea of wrapping their social anxieties in the neat package of a psychological complex appealed to many Americans.

The Inferiority Complex became an all-purpose explanation for problems in many areas of life, ranging from love to parenting to career.

In 1924, \textit{Collier's} ran a story about a woman who was afraid to marry the man she loved for fear that he had an IC and would never amount to anything.

Another popular magazine ran an article called ``Your Child and That Fashionable Complex,'' explaining to moms what could cause in IC in kids and how to prevent or cure one.

\textit{Everyone} had an IC, it seemed; to some it was, paradoxically enough, a mark of distinction.

Lincoln, Napoleon, Teddy Roosevelt, Edison, and Shakespeare - all had suffered from ICs, according to a 1939 \textit{Collier's} article.

``So,'' concluded the magazine, ``if you have a big, husky, in-growing interiority complex you're about as lucky as you could hope to be, provided you have the backbone along with it.''

%
Despite the hopeful tone of this piece, child guidance experts of the 1920s set about helping children to develop winning personalities.

Until then, these professionals had worried mainly about sexually precocious girls and delinquent boys, but now psychologists, social workers, and doctors focused on the everyday child with the ``maladjusted personality'' - particularly shy children.

Shyness could lead to dire outcomes, they warned, from alcoholism to suicide, while an outgoing personality would bring social and financial success.

The experts advised parents to socialize their children well and schools to change their emphasis from book-learning to ``assisting and guiding the developing personality.''

Educators took up this mantle enthusiastically.

By 1950 the slogan of the Mid-Century White House Conference on Children and Youth was ``A healthy personality for every child.''

%
Well-meaning parents of the midcentury agreed that quiet was unacceptable and gregariousness ideal for both girls and boys.

Some discouraged their children from solitary and serious hobbies, like classical music, that could make them unpopular.

They sent their kids to school at increasingly young ages, where the main assignment was learning to socialize.

Introverted children were often singled out as problem cases (a situation familiar to anyone with an introverted child today).

%
William Whyte's \textit{The Organization Man}, a 1956 best-seller, describes how parents and teachers conspired to overhaul the personalities of quiet children.

``Johnny wasn't doing so well at school,'' Whyte recalls a mother telling him.

``The teacher explained to me that he was doing fine on his lessons but that his social adjustment was not as good as it might be.

He would pick just 1 or 2 friends to play with, and sometimes he was happy to remain by himself.''

Parents welcomed such interventions, said Whyte.

``Save for a few odd parents, most are grateful that the schools work so hard to offset tendencies to introversion and other suburban abnormalities.''

%
Parents caught up in this value system were not unkind, or even obtuse; they were only preparing their kids for the ``real world.''

When these children grew older and applied to college and later for their 1st jobs, they faced the same standards of gregariousness.

University admission officers looked not for the most exceptional candidates, but for the most extroverted.

Harvard's provost Paul Buck declared in the late 1940s that Harvard should reject the ``sensitive, neurotic'' type and the ``intellectually over-stimulated'' in favor of boys of the ``healthy extrovert kind.''

In 1950, Yale's president, Alfred Whitney Griswold, declared that the ideal Yalie was not a ``beetle-browed, highly specialized intellectual, but a well-rounded man.''

Another dean told Whyte that ``in screening applications from secondary schools he felt it was only common sense to take into account not only what the college wanted, but what, 4 years later, corporations' recruiters would want.

`They like a pretty gregarious, active type,' he said.

`So we find that the best man is the one who's had an 80 or 85 average in school and plenty of extracurricular activity.

We see little use for the \fbox{``brilliant'' introvert}.'''

%
This college dean grasped very well that the model employee of the midcentury - even one whose job rarely involved dealing with the public, like a research scientist in a corporate lab - was not a deep thinker but a hearty extrovert with a salesman's personality.

``Customarily, whenever the word brilliant is used,'' explains Whyte, ``it either precedes the word `but' (e.g., `We are all for brilliance, but$\ldots$') or is coupled with such words as erratic, eccentric, introvert, screwball, etc.''

``These fellows will be having contact with other people in the organization,'' said one 1950s executive about the hapless scientists in his employ, ``and it helps if they make a good impression.''

%
\fbox{The scientist's job was not only to do the research but also to help sell it}, and that required a hail-fellow-well-met demeanor.

At IBM, a corporation that embodied the ideal of the company man, the sales force gathered each morning to belt out the company anthem, ``Ever Onward,'' and to harmonize on the ``Selling IBM'' song, set to the tune of ``Singin' in the Rain.''

``Selling IBM,'' it began, ``we're selling IBM.

What a glorious feeling the world is our friend.''

The ditty built to a stirring close: ``We're always in trim, we work with a vim.

We're selling, just selling, IBM.''

%
Then they went off to pay their sales calls, proving that the admissions people at Harvard and Yale were probably right: only a certain type of fellow could possibly have been interested in kicking off his mornings this way.

%
The rest of the organization men would have to manage as best they could.

And if the history of pharmaceutical consumption is any indication, many buckled under such pressures.

In 1955 a drug company named Carter-Wallace released the anti-anxiety drug Miltown, reframing anxiety as the natural product of a society that was both dog-eat-dog and relentlessly social.

Miltown was marketed to men and immediately became the fastest-selling pharmaceutical in American history, according to the social historian Andrea Tone.

By 1956 1 of every 20 Americans had tried it; by 1960 a 3rd of all prescriptions from U.S. doctors were for Miltown or a similar drug called Equanil.

``\textsc{anxiety and tension are the commonplace of the age},'' read the Equanil ad.

The 1960s tranquilizer Serentil followed with an ad campaign even more direct in its appeal to improve social performance.

``\textsc{for the anxiety that comes from not fitting in},'' it emphasized.
\begin{center}
    $\star$
\end{center}
Of course, the Extrovert Ideal is not a modern invention.

Extroversion is in our DNA - literally, according to some psychologists.

The trait has been found to be less prevalent in Asia and Africa than in Europe and America, whose populations descend largely from the migrants of the world.

It makes sense, say these researchers, that world travelers were more extroverted than those who stayed home - and that they passed on their traits to their children and their children's children.

``As personality traits are genetically transmitted,'' writes the psychologist Kenneth Olson, ``each succeeding wave of emigrants to a new continent would give rise ovre time to a population of more engaged individuals than reside in the emigrants' continent of origin.''

%
We can also trace our admiration of extroverts to the Greeks, for whom oratory was an exalted skill, and to the Romans, for whom the worst possible punishment was banishment from the city, with its teeming social life.

Similarly, we reverse our founding fathers precisely because they were loudmouths on the subject of freedom: \textit{Give me liberty or give me death!}

Even the Christianity of early American religious revivals, dating back to the 1st Great Awakening of the 18th century, depended on the showmanship of ministers who were considered successful if they caused crowds of normally reserved people to weep and shout and generally lose their decorum.

``Nothing gives me more pain and distress than to see a minister standing almost motionless, coldly plodding on as a mathematician would calculate the distance of the Moon from the Earth,'' complained a religious newspaper in 1837.

%
As this disdain suggests, early Americans revered action and were suspicious of intellect, associating the life of the mind with the languid, ineffectual European aristocracy they had left behind.

The 1828 presidential campaign pitted a former Harvard professor, John Quincy Adams, against Andrew Jackson, a forceful military hero.

A Jackson campaign slogan telling distinguished the 2: ``John Quincy Adams who can write/And Andrew Jackson who can fight.''

%
The victor of that campaign?

The fighter beat the writer, as the cultural historian Neal Gabler puts it.

(John Quincy Adams, incidentally, is considered by political psychologists to be 1 of the few introverts in presidential history.)

%
But the rise of the Culture of Personality intensified such biases, and applied them not only to political and religious leaders, but also to regular people.

And though soap manufacturers may have profited from the new emphasis on charm and charisma, not everyone was pleased with this development.

``Respect for individual human personality has with us reached its lowest point,'' observed 1 intellectual in 1921, ``and it is delightfully ironical that no nation is so constantly talking about personality as we are.

We actually have schools for `self-expression' and `self-development,' although we seem usually to mean the expression and development of the personality of a successful real estate agent.''

%
Another critic bemoaned the slavish attention Americans were starting to pay to entertainers: ``It is remarkable how much attention the stage and things pertaining to it are receiving nowadays from the magazines,'' he grumbled.

Only 20 years earlier - during the Culture of Character, i.e. - such topics would have been considered indecorous; now they had become ``such a large part of the life of society that it has become a topic of conversation among all classes.''

%
Even T. S. Eliot's famous 1915 poem \textit{The Love Song of J. Alfred Prufrock} - in which he laments the need to ``prepare a face to meet the faces that you meet'' - seems a \textit{cri de coeur} about the new demands of self-presentation.

While poets of the previous century had wandered lonely as a cloud through the countryside (Wordsworth, in 1802) or repaired in solitude to Walden Pond (Thoreau, in 1845), Eliot's Prufrock mostly worries about being looked at by ``eyes that fix you in a formulated phrase'' and pin you, wriggling, to a wall.
\begin{center}
    $\star$
\end{center}
Fast-forward nearly 100 years, and Prufrock's protest is enshrined in high school syllabi, where it's dutifully memorized, then quickly forgotten, by teens increasingly skilled at shaping their own online and offline personae.

These students inhabit a world in which status, income, and self-esteem depend more than ever on the ability to meet the demands of the Culture of Personality.

The pressure to entertain, to sell ourselves, and never to be visibly anxious keeps ratcheting up.

The number of Americans who considered themselves shy increased from 40\% in the 1970s to 50\% in the 1990s, probably because we measured ourselves against ever higher standards of fearless self-presentation.

``Social anxiety disorder'' - which essentially means pathological shyness - is now thought to afflict nearly 1 in 5 of us.

The most recent version of the \textit{Diagnostic and Statistical Manual (DSM-IV)}, the psychiatrist's bible of mental disorders, considers the fear of public speaking to be a pathology - not an annoyance, not a disadvantage, but a \textit{disease} - if it interferes with the sufferer's job performance.

``It's not enough,'' 1 senior manager at Eastman Kodak told the author Daniel Goleman, ``to be able to sit at your computer excited about a fantastic regression analysis if you're squeamish about presenting those results to an executive group.''

(Apparently it's OK to be squeamish about doing a regression analysis if you're excited about giving speeches.)

%
But perhaps the best way to take the measure of the 21st-century Culture of Personality is to return to the self-help arena.

Today, a full century after Dale Carnegie launched that 1st public-speaking workshop at the YMCA, his best-selling book \textit{How to Win Friends and Influence People} is a staple of airport bookshelves and business best-seller lists.

The Dale Carnegie Institute still offers updated versions of Carnegie's original classes, and the ability to communicate fluidly remains a core feature of the curriculum.

Toastmasters, the nonprofit organization established in 1924 whose members meet weekly to practice public speaking and whose founder declared that ``all talking is selling and all selling involves talking,'' is still thriving, with more than 12,500 chapters in 113 countries.

%
The promotional video on Toastmaster's website features a skit in which 2 colleagues, Eduardo and Sheila, sit in the audience at the ``6th Annual Global Business Conference'' as a nervous speaker stumbles through a pitiful presentation.

%
``I'm so glad I'm not him,'' whispers Eduardo.

``You're joking, right?'' replies Sheila with a satisfied smile.

``Don't you remember last month's sales presentation to those new clients?

I thought you were going to faint.''

%
``I wasn't that bad, was I?''

%
``Oh, you were that bad.

Really bad.

Worse, even.''

%
Eduardo looks suitably ashamed, while the rather insensitive Sheila seems oblivious.

%
``But,'' says Sheila, ``you can fix it.

You can do better$\ldots$

Have you ever heard of Toastmasters?''

%
Sheila, a young and attractive brunette, hauls Eduardo to a Toastmasters meeting.

There she volunteers to perform an exercise called ``Truth or Lie,'' in which she's supposed to tell the group fo 15-odd participants a story about her life, after which they decide whether or not to believe her.

%
``I bet I can fool everyone,'' she whispers to Eduardo sotto voce as she marches to the podium.

She spins an elaborate tale about her years as an opera singer, concluding with her poignant decision to give it all up to spend more time with her family.

When she's finished, the toastmaster of the evening asks the group whether they believe Sheila's story.

All hands in the room go up.

The toastmaster turns to Sheila and asks whether it was true.

%
``I can't even carry a tune!'' she beams triumphantly.

%
Sheila comes across as disingenuous, but also oddly sympathetic.

Like the anxious readers of the 1920s personality guides, she's only trying to get ahead at the office.

``There's no much competition in my work environment,'' she confides to the camera, ``that it makes it more important than ever to keep my skills sharp.''

%
But what do ``sharp skills'' look like?

Should we become so proficient at self-presentation that we can dissemble without anyone suspecting?

Must we learn to stage-manage our voices, gestures, and body language until we can tell - sell - any story we want?

These seem venal aspirations, a marker of how far we've come - and not in a good way - since the days of Dale Carnegie's childhood.

%
Dale's parents had high moral standards; they wanted their son to pursue a career in religion or education, not sales.

\fbox{It seems unlikely that they would have approved of a self-improvement technique called ``Truth or Lie.''}

Or, for that matter, of Carnegie's best-selling advice on how to get people to admire you and do your bidding.

\textit{How to Win Friends and Influence People} is full of chapter titles like ``Making People Glad to Do What You Want'' and ``How to Make People Like You Instantly.''

%
All of which raises the question, \textit{how did we go from Character to Personality without realizing that we had sacrificed something meaningful along the way?}

%------------------------------------------------------------------------------%

\section{\textsc{The Myth of Charismatic Leadership}: The Culture of Personality, a Hundred Years Later}

%------------------------------------------------------------------------------%

\section{\textsc{When Collaboration Kills Creativity}: The Rise of the New Groupthink \& the Power of Working Alone}

%------------------------------------------------------------------------------%

\part{\textsc{Your Biology, Your Self?}}

\section{\textsc{Is Temperament Destiny?}: Nature, Nurture, \& the Orchid Hypothesis}


%------------------------------------------------------------------------------%

\section{\textsc{Beyond Temperament}: The Role of Free Will (\& the Secret of Public Speaking for Introverts)}

%------------------------------------------------------------------------------%

\section{\textsc{``Franklin Was a Politician, But Eleanor Spoke out of Conscience''}: Why Cool is Overrated?}

%------------------------------------------------------------------------------%

\section{\textsc{Why Did Wall Street Crash \& Warren Buffett Prosper?}: How Introverts \& Extroverts Think (\& Process Dopamine) Differently}

%------------------------------------------------------------------------------%

\part{\textsc{Do All Cultures Have an Extrovert Ideal?}}

\section{\textsc{Soft Power}: Asian-Americans \& the Extrovert Ideal}

%------------------------------------------------------------------------------%

\part{\textsc{How to Love, How to Work}}

\section{\textsc{When Should You Act More Extroverted Than You Really Are?}}

%------------------------------------------------------------------------------%

\section{\textsc{The Communication Gap}: How to Talk to Members of the Opposite Type}

%------------------------------------------------------------------------------%

\section{\textsc{On Cobblers \& Generals}: How to Cultivate Quiet Kids in a World That Can't Hear Them}

%------------------------------------------------------------------------------%

\section{\textsc{Conclusion}: Wonderland}

%------------------------------------------------------------------------------%

\section{A Note on the Dedication}

%------------------------------------------------------------------------------%

\section{A Note on the Words \textit{Introvert} \& \textit{Extrovert}}










%------------------------------------------------------------------------------%

\printbibliography[heading=bibintoc]
\end{document}